\documentclass[11pt]{article}
\usepackage{amsmath,amssymb,amsthm}
\newtheorem{proposition}{Proposition}
\title{A small adaptive-pole test with fixed rank masks}
\author{}\date{}
\begin{document}\maketitle
\paragraph{Status and scope.}
This note gives a complete written finite-dimensional construction and an
exact rational checker. It is not Lean-checked and gives no growing
query-to-head ratio. The poles depend on the desired signs. Its purpose
is to test a different mechanism from a fixed token pool with a
butterfly permutation, after the common-pool compatibility obstructions.

Let $t_0<t_1<t_2<t_3$ be fixed positive queries. A root uses eight
ordered positive poles, positive residues $c_j$, and the alternating
rank mask $(-1)^j$. Its two descendant masks are $(-1)^j$ on the
even-rank and odd-rank subsequences respectively. Each descendant is
sampled at $t_0,t_3$ and may have its own scalar bias.
A global exchange of the two root children reverses the root mask and
exchanges the two descendant labels.

\begin{proposition}
Every four-bit root word and both independently specified two-bit
child words are realized by this system, after allowing the one global
exchange. The poles can lie in an arbitrarily small positive interval.
Moreover four fixed ordered pole shapes, scaled by one sufficiently
small common positive parameter, suffice for every word.
\end{proposition}
\begin{proof}
For a desired root word $\sigma$, put one root at the midpoint of
$(t_j,t_{j+1})$ exactly when $\sigma_j\ne\sigma_{j+1}$. Let $N$ be
the product of those factors, with sign normalized to be negative on
$(0,t_0)$. Thus $\deg N\leq3$, its query word is $-\sigma_0\sigma$,
and it is nonzero at zero. For ordered poles $p_j$, set
\[
 Q(t)=\prod_{j=0}^7(t-p_j),\qquad
 c_j=\left|\frac{N(p_j)}{Q'(p_j)}\right|.
\]
Because $\operatorname{sign}Q'(p_j)=-(-1)^j$, partial fractions give
\[
 \sum_{j=0}^7\frac{(-1)^j c_j}{t-p_j}=\frac{N(t)}{Q(t)}.
\]
The optional global exchange therefore supplies exactly $\sigma$.

Write $p_j=\epsilon x_j$ and $P(x)=\prod_j(x-x_j)$. The two
unbiased child scores $G_b$, $b=0,1$, satisfy
\[
 \epsilon^7\bigl(G_b(t_0)-G_b(t_3)\bigr)
 \longrightarrow
 |N(0)|\left(\frac1{t_0}-\frac1{t_3}\right)L_b,
 \qquad
 L_b=\sum_{j=0}^3\frac{(-1)^j}{|P'(x_{2j+b})|}.
\]
Choose the pole shape from the following table:
\[
\begin{array}{c|rrrrrrrr}
 (\operatorname{sign}L_0,\operatorname{sign}L_1)&\multicolumn{8}{c}{(x_0,\ldots,x_7)}\\
 (+,+)&98&166&237&303&368&387&486&538\\
 (-,-)&98&180&270&332&340&419&434&441\\
 (+,-)&38&73&120&150&154&201&300&393\\
 (-,+)&25&119&206&287&375&448&474&573
\end{array}
\]
The indicated signs are exact rational inequalities, verified by the
companion checker. Hence every pair of child trends is available for
all sufficiently small positive $\epsilon$.

A child word with opposite signs needs the trend whose sign is its
first label; the bias $-(G_b(t_0)+G_b(t_3))/2$ realizes it. A child
word with equal labels is obtained by a sufficiently large bias of
that sign and puts no condition on the trend. If the root children
were exchanged, exchange their required trends before selecting the
table row. There are only finitely many root words and pole shapes,
so their strict limiting inequalities hold at one common sufficiently
small positive scale. This scale can be reduced further to put every
pole inside any prescribed positive neighborhood of zero.
\end{proof}

\paragraph{Verification.}
\texttt{fixed-rank-moving-poles-check.py} checks the four leading-sign
pairs exactly. At $t_j=1+j/100$ and $\epsilon=10^{-5}$ it also checks
all $16$ root words for all four shapes, including the exact root
partial-fraction identity and both strict child trends.

\paragraph{What does not follow.}
At the next size, the two children have four queries and need more
than a single trend. Generic compact pole shapes make their leading
scores constant multiples of $1/t$, so adding a bias does not produce
arbitrary four-query words. One would need target-adaptive moment
cancellations or another degeneracy. The finite construction above
provides neither that invariant nor a superconstant saving.

\paragraph{A partition caution.}
If $Q=R^2-\delta^2$, the two polynomial-preimage groups $R=\delta$
and $R=-\delta$ are generally not the even-rank and odd-rank groups:
the order inside each close pair reverses when $R'$ changes sign.
With independently chosen masks one can still use those groups.
Their alternating masked transforms are scalar multiples of the
proper parts of $N/(R\mp\delta)$. With a fixed pole hierarchy these
are all linear functions of the coefficients of $N$, including any
fixed-profile corrections whose coefficients depend linearly on $N$.
That identity alone does not supply more independent signs than the
numerator coefficient dimension.
\end{document}
